%===============================================================================
% БАКАЛАВРЫН ТӨГСӨЛТИЙН АЖЛЫН ИЛТГЭЛ — V3 (Element бүрийг засварлах боломжтой)
% М.Түвшинжаргал (B221940619), Удирдагч: Доктор Д.Эрдэнэтуяа
%
% ЗАСВАРЛАХ ЗААВАР:
%   • Element бүр өөрийн мөр, тайлбартай.
%   • Position constants (\def-ээр) frame бүрийн дээр байрлана — координат
%     өөрчлөхөд тэр макрог л солино.
%   • Текстийг node-ын { ... } доторх хэсэгт засна.
%   • Өнгийг \definecolor блокт өөрчилнө (preamble хэсэгт).
%   • Шинэ element нэмэхдээ "% --- ШИНЭ ELEMENT ---" гэж тэмдэглэн нэмж болно.
%===============================================================================

\documentclass[12pt,aspectratio=169]{beamer}

%--- Encoding & fonts ----------------------------------------------------------
\usepackage[T2A]{fontenc}
\usepackage[utf8]{inputenc}
\usepackage[mongolian,english]{babel}

%--- Packages ------------------------------------------------------------------
\usepackage{tikz}
\usetikzlibrary{arrows.meta,shapes.geometric,shapes.symbols,positioning,calc,
                decorations.pathmorphing,decorations.markings,fit,backgrounds}
\usepackage{pgfplots}
\pgfplotsset{compat=1.18}
\usepackage{xcolor}
\usepackage{array,colortbl}
\usepackage{amsmath,amssymb}
\usepackage{tabularx}
\usepackage{multirow}

%===============================================================================
% ӨНГӨН СХЕМ — ӨӨРЧЛӨХДӨӨ ЭНД ЗАСНА
%===============================================================================
\definecolor{ink}{RGB}{30,30,40}            % бараг хар бэх
\definecolor{paper}{RGB}{253,251,243}       % цайвар цаасны өнгө
\definecolor{paper2}{RGB}{246,242,228}      % бага зэрэг хар цаас
\definecolor{accent}{RGB}{200,80,40}        % улаан тэмдэглэгээ
\definecolor{accent2}{RGB}{60,120,80}       % ногоон тэмдэглэгээ
\definecolor{accent3}{RGB}{50,90,160}       % хөх тэмдэглэгээ
\definecolor{rowgray}{RGB}{248,245,235}
\definecolor{titlebg}{RGB}{40,55,90}
\definecolor{titlefg}{RGB}{250,248,238}

%===============================================================================
% TIKZ STYLE-УУД — БҮХ ДИАГРАМЫН СУУРЬ СТИЛЬ
% hd        — гар-зурсан төстэй (rounded, line cap)
% hdbox=#1  — суурь box (#1 = өнгө)
% hdarrow   — гар-зурсан сум
% hddash    — тасархай шугам
% hdellipse — Use Case оваль
% hddiamond — шийдвэрийн diamond
%===============================================================================
\tikzset{every node/.style={align=center}}

\tikzset{
  hd/.style       = {line cap=round, line join=round, line width=0.7pt},
  hdline/.style   = {decoration={random steps, segment length=6pt, amplitude=0.4pt},
                     decorate, line cap=round, line join=round, line width=0.6pt},
  hdbox/.style    = {hd, draw=ink, fill=#1, rounded corners=3pt, inner sep=4pt, font=\small},
  hdbox/.default  = paper,
  hdarrow/.style  = {hdline, -{Latex[length=2.2mm,width=2mm]}, draw=ink},
  hddash/.style   = {hdline, draw=ink, dashed},
  hdellipse/.style= {hd, draw=ink, fill=#1, ellipse, inner sep=3pt, font=\small},
  hdellipse/.default = paper,
  hddiamond/.style= {hd, draw=ink, fill=#1, diamond, aspect=2.4, inner sep=2pt, font=\scriptsize},
  hddiamond/.default = paper2,
}

%--- Beamer theme --------------------------------------------------------------
\usetheme{default}
\setbeamercolor{frametitle}{bg=titlebg,fg=titlefg}
\setbeamercolor{title}{fg=titlebg}
\setbeamercolor{author}{fg=titlebg}
\setbeamercolor{date}{fg=titlebg!70}
\setbeamercolor{structure}{fg=titlebg}
\setbeamercolor{itemize item}{fg=accent}
\setbeamercolor{itemize subitem}{fg=accent3}
\setbeamercolor{normal text}{fg=ink}
\setbeamerfont{frametitle}{size=\large,series=\bfseries}
\setbeamertemplate{navigation symbols}{}
\setbeamertemplate{frametitle}{%
  \nointerlineskip
  \begin{beamercolorbox}[wd=\paperwidth,ht=2.6ex,dp=0.8ex,leftskip=1em,rightskip=1em]{frametitle}
    \usebeamerfont{frametitle}\insertframetitle
  \end{beamercolorbox}%
  \nointerlineskip
  {\color{accent}\hrule height 1.2pt}%
}
%===============================================================================
% FOOTER — явцын зурвас + бүлгийн дугаар + хуудасны дугаар
%   ① Явцын зурвас (accent өнгө, бүрэн өргөн рүү ургана)
%   ② Хуудасны дугаар  N / T
%   ③ Бүлгийн цэгүүд  ● ● ○ ○ ○  (5 хэсэг)
%   ④ Одоогийн бүлгийн нэр  § Бүлэг X
%   ⑤ Зохиогч ба гарчиг
%===============================================================================
\newlength{\beamerprogress}
\setbeamercolor{footprog}{bg=paper,fg=paper}
\setbeamercolor{footinfo}{bg=paper,fg=ink}

\setbeamertemplate{footline}{%
  \leavevmode%
  %── ① Явцын зурвас ──────────────────────────────────────────────
  \pgfmathsetmacro{\fp}{\insertframenumber/\inserttotalframenumber}%
  \setlength{\beamerprogress}{\fp\paperwidth}%
  \begin{beamercolorbox}[wd=\paperwidth,ht=2.2pt,dp=0pt]{footprog}%
    {\color{ink!15}\rule{\paperwidth}{2.2pt}}%
    \hskip-\paperwidth%
    {\color{accent}\rule{\beamerprogress}{2.2pt}}%
  \end{beamercolorbox}%
  %── ②③④⑤ Хуудасны мэдээлэл ─────────────────────────────────────
  \begin{beamercolorbox}[wd=\paperwidth,ht=2.6ex,dp=0.9ex,%
                         leftskip=0.55em,rightskip=0.55em]{footinfo}%
    \tiny%
    % ② Хуудасны дугаар
    \textbf{\color{accent}\insertframenumber}%
    {\color{ink!55}\ /\ \inserttotalframenumber}%
    \hspace*{0.5em}%
    % ③ Бүлгийн дугаарын цэгүүд (бүх хэсэг = 5)
    \foreach \si in {1,...,5}{%
      \ifnum\si=\value{section}%
        \textcolor{accent}{\textbullet}%          одоогийн бүлэг
      \else\ifnum\si<\value{section}%
        \textcolor{ink!62}{\textbullet}%          өнгөрсөн бүлэг
      \else%
        \textcolor{ink!20}{\textbullet}%          ирэх бүлэг
      \fi\fi%
      \hspace{0.1em}%
    }%
    % ④ Одоогийн бүлгийн нэр
    \ifnum\value{section}>0%
      \hspace*{0.38em}%
      {\color{accent3}\bfseries\S\,\insertsectionhead}%
    \fi%
    \hfill%
    % ⑤ Зохиогч
    М.Түвшинжаргал --- Үүлэн орчинд суурилсан спорт заалны цаг захиалгын систем%
  \end{beamercolorbox}%
}

\title{Үүлэн орчинд суурилсан спорт заалны цаг захиалгын\\3 давхаргат системийн загварчлал ба хэрэгжилт}
\author{М.Түвшинжаргал (B221940619)}
\institute{Өмнөговь Технологийн Дээд Сургууль}
\date{2026}

\begin{document}

%###############################################################################
% SLIDE 1 — НҮҮР
%###############################################################################
\begin{frame}[plain]
  % --- Дэвсгэр ба хүрээ -----------------------------------------------------
  \begin{tikzpicture}[remember picture, overlay]
    \fill[paper] (current page.north west) rectangle (current page.south east);
    \draw[draw=accent, line width=1.2pt, rounded corners=4pt]
      ($(current page.north west)+(1.2,-0.8)$) rectangle ($(current page.south east)+(-1.2,1.2)$);
  \end{tikzpicture}

  \vspace{1.2cm}
  \begin{center}
    % --- Сургуулийн нэр -----------------------------------------------------
    {\small\bfseries\color{titlebg} ӨМНӨГОВЬ ТЕХНОЛОГИЙН ДЭЭД СУРГУУЛЬ}\\[2pt]
    {\footnotesize\color{ink} Мэдээллийн технологийн тэнхим}\\[2pt]
    {\footnotesize\color{ink}\itshape Бакалаврын төгсөлтийн ажлын илтгэл}\\[1.2cm]

    % --- Гарчгийн bлок ----------------------------------------------------
    \begin{tikzpicture}
      \node[hdbox=paper2, line width=1pt, text width=12cm,
            font=\large\bfseries\color{titlebg}, inner sep=14pt] {
        ҮҮЛЭН ОРЧИНД СУУРИЛСАН СПОРТ ЗААЛНЫ\\
        ЦАГ ЗАХИАЛГЫН 3 ДАВХАРГАТ СИСТЕМИЙН\\
        ЗАГВАРЧЛАЛ БА ХЭРЭГЖИЛТ
      };
    \end{tikzpicture}\\[1.2cm]

    % --- Илтгэгч / Удирдагч --------------------------------------------------
    \begin{tabular}{rl}
      \textbf{Илтгэгч:}  & М.Түвшинжаргал (B221940619)\\[3pt]
      \textbf{Удирдагч:} & Доктор (Ph.D) Д.Эрдэнэтуяа\\
    \end{tabular}\\[1cm]
    {\small\color{ink!70} Улаанбаатар хот, 2026 он}
  \end{center}
\end{frame}

%###############################################################################
% SLIDE 2 — АГУУЛГА
% Element-ууд: 5 хэсгийн картууд (a1..a5)
%###############################################################################
\begin{frame}{Агуулга}
  \vspace{0.1cm}
  \centering
  \begin{tikzpicture}
    % Row 1 — Үндэслэл  (center y = 2.30)
    \filldraw[fill=paper,  draw=ink!15, line width=0.3pt] (-5.3,1.80) rectangle (5.3,2.80);
    \fill[ink]  (-5.3,1.80) rectangle (-4.70,2.80);
    \node[font=\bfseries\small, text=paper]                   at (-5.00,2.30) {1};
    \node[anchor=west, font=\bfseries\small,  text=ink]       at (-4.50,2.46) {Үндэслэл};
    \node[anchor=west, font=\scriptsize\itshape, text=ink!60] at (-4.50,2.14) {Асуудал · зорилго · шийдэл};
    % Row 2 — Бүлэг 1  (center y = 1.15)
    \filldraw[fill=paper,  draw=ink!15, line width=0.3pt] (-5.3,0.65) rectangle (5.3,1.65);
    \fill[ink]  (-5.3,0.65) rectangle (-4.70,1.65);
    \node[font=\bfseries\small, text=paper]                   at (-5.00,1.15) {2};
    \node[anchor=west, font=\bfseries\small,  text=ink]       at (-4.50,1.31) {Бүлэг 1: Онол};
    \node[anchor=west, font=\scriptsize\itshape, text=ink!60] at (-4.50,0.99) {3 давхаргат архитектур · технологийн стек};
    % Row 3 — Бүлэг 2  (center y = 0.00)
    \filldraw[fill=paper,  draw=ink!15, line width=0.3pt] (-5.3,-0.50) rectangle (5.3,0.50);
    \fill[ink]  (-5.3,-0.50) rectangle (-4.70,0.50);
    \node[font=\bfseries\small, text=paper]                   at (-5.00,0.00) {3};
    \node[anchor=west, font=\bfseries\small,  text=ink]       at (-4.50,0.16) {Бүлэг 2: Хөгжүүлэлт};
    \node[anchor=west, font=\scriptsize\itshape, text=ink!60] at (-4.50,-0.16) {Use case · UI/UX · AI чатбот · өгөгдлийн бүтэц};
    % Row 4 — Бүлэг 3  (center y = -1.15)
    \filldraw[fill=paper,  draw=ink!15, line width=0.3pt] (-5.3,-1.65) rectangle (5.3,-0.65);
    \fill[ink]  (-5.3,-1.65) rectangle (-4.70,-0.65);
    \node[font=\bfseries\small, text=paper]                   at (-5.00,-1.15) {4};
    \node[anchor=west, font=\bfseries\small,  text=ink]       at (-4.50,-0.99) {Бүлэг 3: Үүлэн орчин};
    \node[anchor=west, font=\scriptsize\itshape, text=ink!60] at (-4.50,-1.31) {Docker · AWS EC2 · CI/CD · туршилт};
    % Row 5 — Дүгнэлт  (center y = -2.30, paper2 bg)
    \filldraw[fill=paper2, draw=ink!15, line width=0.3pt] (-5.3,-2.80) rectangle (5.3,-1.80);
    \fill[ink!75](-5.3,-2.80) rectangle (-4.70,-1.80);
    \node[font=\bfseries\small, text=paper]                   at (-5.00,-2.30) {5};
    \node[anchor=west, font=\bfseries\small,  text=ink]       at (-4.50,-2.14) {Дүгнэлт};
    \node[anchor=west, font=\scriptsize\itshape, text=ink!60] at (-4.50,-2.46) {Үр дүн · ач холбогдол · цаашдын хөгжил};
  \end{tikzpicture}
\end{frame}

%###############################################################################
% §1 — ҮНДЭСЛЭЛ (хэсэг 1)
\section[Үндэслэл]{Үндэслэл}
%###############################################################################
% SLIDE 3 — ҮНДЭСЛЭЛ (Асуудал ↔ Шийдэл)
% Element-ууд:
%   h1, h2  — толгой headers
%   p1, p2  — асуудал/шийдлийн жагсаалт
%   k1..k3  — доод 3 ключевой box
%   ar1..ar3— холбох сумнууд
%###############################################################################
\begin{frame}{Үндэслэл --- Асуудал ба шийдэл}
  \vspace{0.2cm}
  \centering
  \begin{tikzpicture}[font=\small]
    % --- Position constants -------------------------------------------------
    \def\hLY{4}            % толгойн y
    \def\pLY{3.4}          % paragraph-ийн y (north)
    \def\kY{-2.2}          % доод box-ын y
    \def\colLeftX{-4}      % зүүн баганын x
    \def\colRightX{4}      % баруун баганын x

    % --- Element: Асуудлууд толгой ------------------------------------------
    \node[hdbox=paper2, font=\bfseries, minimum width=6cm] (h1)
      at (\colLeftX,\hLY) {Тулгамдаж буй асуудлууд};
    % --- Element: Асуудлуудын жагсаалт --------------------------------------
    \node[hdbox=paper, text width=5.5cm, align=left, anchor=north] (p1)
      at (\colLeftX,\pLY) {%
        \textbullet\ Утсаар захиалахад шугам завгүй байх\\[2pt]
        \textbullet\ Биечлэн очиж захиалах хүндрэл\\[2pt]
        \textbullet\ Цагийн мэдээлэл ил тод бус\\[2pt]
        \textbullet\ Захиалгаа цуцлах боломжгүй\\[2pt]
        \textbullet\ Сануулга авах боломжгүй};

    % --- Element: Шийдэл толгой ----------------------------------------------
    \node[hdbox=paper2, font=\bfseries, minimum width=6cm] (h2)
      at (\colRightX,\hLY) {Системийн шийдэл};
    % --- Element: Шийдлийн жагсаалт ------------------------------------------
    \node[hdbox=paper, text width=5.5cm, align=left, anchor=north] (p2)
      at (\colRightX,\pLY) {%
        \textbullet\ \textbf{Мобайл аппликейшн} (Flutter)\\[2pt]
        \textbullet\ Бодит цагийн хуваарь (Firestore)\\[2pt]
        \textbullet\ DBK-xxxxx захиалгын код\\[2pt]
        \textbullet\ Локал мэдэгдэл (1 цагийн өмнө)\\[2pt]
        \textbullet\ AI чатбот туслах};

    % --- Element: Доод box 1 — Мобайл аппликейшн ---------------------------
    \node[hdbox=paper2, font=\bfseries\color{accent3}, minimum width=4.2cm] (k1)
      at (-4.5,\kY) {Мобайл аппликейшн};
    % --- Element: Доод box 2 — Архитектур ---------------------------------
    \node[hdbox=paper2, font=\bfseries\color{accent3}, minimum width=4.2cm] (k2)
      at (0,\kY) {3 давхаргат архитектур};
    % --- Element: Доод box 3 — Үүлэн дэд бүтэц -----------------------------
    \node[hdbox=paper2, font=\bfseries\color{accent3}, minimum width=4.2cm] (k3)
      at (4.5,\kY) {Үүлэн дэд бүтэц};

    % --- Element: Сум 1 (асуудал → мобайл) -----------------------------------
    \draw[hdarrow] (p1.south) to[bend right=20] (k1.north);
    % --- Element: Сум 2 (шийдэл → үүлэн) -------------------------------------
    \draw[hdarrow] (p2.south) to[bend left=20]  (k3.north);
    % --- Element: Сум 3 (төв архитектур руу) ---------------------------------
    \draw[hdarrow, draw=accent] ($(p1.south east)+(0.4,0)$) to[bend right=10] (k2.north);
  \end{tikzpicture}
\end{frame}

%###############################################################################
% SLIDE 4 — ЗОРИЛГО (6 sub-card)
% Element-ууд:
%   topBox       — дээд талын зорилгын тодорхойлолт
%   badgeZorigo  — "ЗОРИЛГО" tag
%   card-1..6    — 6 ширхэг карт (icon+title+desc)
%###############################################################################
\begin{frame}{Зорилго}
  \vspace{-0.05cm}
  % Icon ID-ний макро (\ifx харьцуулалтанд хэрэгтэй)
  \def\iconlock{lock}\def\iconphone{phone}\def\iconai{ai}%
  \def\iconclock{clock}\def\iconbell{bell}\def\iconcourt{court}%
  \begin{tikzpicture}[font=\footnotesize]

    % --- Element: Зорилгын топ-box -------------------------------------------
    \node[draw=accent, line width=0.8pt, fill=paper2, rounded corners=4pt,
          text width=14cm, font=\small\itshape, inner sep=8pt] (topBox)
      at (0,3.0) {%
        Даланзадгад хотын спортын заалны захиалгыг хялбар, хурдан, найдвартай болгох мобайл аппликейшн хөгжүүлж, хэрэглэгч болон заалны эздэд ашиглуулах.
      };
    % --- Element: ЗОРИЛГО тэмдэг ---------------------------------------------
    \node[fill=accent, text=white, font=\tiny\bfseries, inner sep=3pt,
          rounded corners=2pt] (badgeZorilgo) at (-5.7,3.7) {ЗОРИЛГО};

    % --- 6 картын тохиргоо --------------------------------------------------
    \def\cw{4.6}    % картын өргөн (см)
    \def\ch{2.0}    % картын өндөр (см)
    % iconID/i/j/colour/title/description гэсэн дарааллаар бичсэн.
    \foreach \icon/\i/\j/\col/\ttl/\dsc in {%
      lock/0/0/accent/{Аюулгүй, хялбар захиалга}/{Хэдхэн товшилтоор спорт заал захиалах},%
      phone/1/0/accent3/{Мобайл аппликейшн}/{Flutter --- iOS \& Android},%
      ai/2/0/accent2/{AI Chat туслах}/{Groq Llama 3.1, лавлагаа өгөх},%
      clock/0/1/accent3/{Бодит цагийн мэдээлэл}/{Firebase Firestore real-time},%
      bell/1/1/accent/{Мэдэгдлийн систем}/{flutter\_local\_notifications},%
      court/2/1/accent2/{Хагас / бүтэн талбай}/{Нэг slot-д хоёр захиалга боломж}}{

        \pgfmathsetmacro{\xpos}{-7.0+\i*4.85}
        \pgfmathsetmacro{\ypos}{0.4-\j*2.25}

        % Element: сүүдэр
        \node[draw=ink!30, fill=ink!10, rounded corners=4pt,
              minimum width=\cw cm, minimum height=\ch cm, anchor=north west]
          at ({\xpos+0.08},{\ypos-0.08}) {};
        % Element: үндсэн карт
        \node[draw=\col, line width=0.9pt, fill=paper, rounded corners=4pt,
              minimum width=\cw cm, minimum height=\ch cm, anchor=north west,
              inner sep=0pt] (card\i\j) at (\xpos,\ypos) {};
        % Element: зүүн талын өнгөт зурвас
        \fill[\col] ($(card\i\j.north west)+(0,0)$) rectangle
                    ($(card\i\j.south west)+(0.18,0)$);
        % Element: дугаар бөмбөлөг (1-6)
        \pgfmathtruncatemacro{\num}{\j*3+\i+1}
        \node[draw=\col, fill=\col, text=white, line width=0.8pt,
              circle, inner sep=0pt, minimum size=0.55cm,
              font=\bfseries\scriptsize, anchor=center]
          at ($(card\i\j.north east)+(-0.4,-0.4)$) {\num};
        % Element: иконы (TikZ-ээр зурсан)
        \coordinate (ic) at ($(card\i\j.north west)+(0.55,-0.55)$);
        \begin{scope}[shift={(ic)}]
          \ifx\icon\iconlock
            \draw[\col, line width=0.6pt] (-0.2,-0.05) rectangle (0.2,-0.35);
            \draw[\col, line width=0.6pt] (-0.13,-0.05) arc (180:0:0.13 and 0.18);
          \else\ifx\icon\iconphone
            \draw[\col, line width=0.6pt, rounded corners=1.5pt] (-0.13,-0.4) rectangle (0.13,0.0);
            \draw[\col, line width=0.4pt] (-0.05,-0.36) -- (0.05,-0.36);
            \draw[\col, line width=0.4pt] (-0.08,-0.04) -- (0.08,-0.04);
          \else\ifx\icon\iconai
            \draw[\col, line width=0.6pt, rounded corners=1pt] (-0.2,-0.35) rectangle (0.2,-0.05);
            \draw[\col, line width=0.5pt] (0,-0.05) -- (0,0.05);
            \fill[\col] (0,0.07) circle (0.04);
            \fill[\col] (-0.1,-0.18) circle (0.03);
            \fill[\col] (0.1,-0.18) circle (0.03);
            \draw[\col, line width=0.4pt] (-0.07,-0.28) -- (0.07,-0.28);
          \else\ifx\icon\iconclock
            \draw[\col, line width=0.6pt] (0,-0.2) circle (0.2);
            \draw[\col, line width=0.6pt] (0,-0.2) -- (0,-0.07);
            \draw[\col, line width=0.6pt] (0,-0.2) -- (0.1,-0.2);
            \draw[\col, line width=0.4pt] (0,-0.42) -- (0,-0.45);
            \draw[\col, line width=0.4pt] (0.22,-0.2) -- (0.25,-0.2);
          \else\ifx\icon\iconbell
            \draw[\col, line width=0.6pt] (-0.2,-0.32) -- (-0.2,-0.18)
              arc (180:0:0.2 and 0.22) -- (0.2,-0.32) -- cycle;
            \draw[\col, line width=0.4pt] (-0.06,-0.4) arc (180:360:0.06);
            \draw[\col, line width=0.4pt] (-0.05,0.04) -- (0.05,0.04);
          \else\ifx\icon\iconcourt
            \draw[\col, line width=0.6pt] (-0.25,-0.35) rectangle (0.25,-0.05);
            \draw[\col, line width=0.5pt] (0,-0.35) -- (0,-0.05);
            \draw[\col, line width=0.4pt] (0,-0.2) circle (0.06);
          \fi\fi\fi\fi\fi\fi
        \end{scope}
        % Element: гарчиг
        \node[font=\scriptsize\bfseries, text=\col, anchor=north west,
              text width=\cw cm-2.0cm, align=left]
          at ($(card\i\j.north west)+(0.95,-0.25)$) {\ttl};
        % Element: тайлбар
        \node[font=\tiny\color{ink!85}, anchor=south west,
              text width=\cw cm-0.6cm, align=left]
          at ($(card\i\j.south west)+(0.35,0.15)$) {\dsc};
      }

  \end{tikzpicture}
\end{frame}

%###############################################################################
% §2 — БҮЛЭГ 1 (хэсэг 2)
\section[Бүлэг 1]{Бүлэг 1: 3 давхаргат архитектур}
%###############################################################################
% SLIDE 5 — БҮЛЭГ 1: 3 ДАВХАРГАТ АРХИТЕКТУРЫН ОНОЛ
% Element-ууд:
%   pres, app, data — 3 давхар layer
%   arrPA, arrAD     — давхрууд хооронд сум
%   labelL, labelM, labelR — хажуугийн товч тэмдэг
%###############################################################################
\begin{frame}{Бүлэг 1: 3 давхаргат архитектурын онол}
  \centering
  \begin{tikzpicture}[font=\small]
    % --- Position constants -------------------------------------------------
    \def\layerW{12}     % давхрын өргөн (cm)
    \def\layerH{1.4}    % давхрын өндөр (cm)
    \def\presY{3}       % presentation tier y
    \def\appY{0.4}      % application tier y
    \def\dataY{-2.2}    % data tier y

    % --- Element: Presentation Tier (Flutter клиент) ------------------------
    \node[hdbox=paper, minimum width=\layerW cm, minimum height=\layerH cm,
          text width=11.5cm, font=\bfseries\color{accent3}] (pres)
      at (0,\presY) {Presentation Tier --- Flutter Мобайл клиент\\
        {\footnotesize\mdseries\color{ink} (Android, iOS) UI/UX дэлгэц, навигаци, локал мэдэгдэл}};

    % --- Element: Application Tier (Node.js бэкэнд) -------------------------
    \node[hdbox=paper, minimum width=\layerW cm, minimum height=\layerH cm,
          text width=11.5cm, font=\bfseries\color{accent2}] (app)
      at (0,\appY) {Application Tier --- Node.js / Express REST API\\
        {\footnotesize\mdseries\color{ink} (AWS EC2, Docker контейнер) бизнесийн логик, email баталгаажуулалт}};

    % --- Element: Data Tier (Firebase) ---------------------------------------
    \node[hdbox=paper, minimum width=\layerW cm, minimum height=\layerH cm,
          text width=11.5cm, font=\bfseries\color{accent}] (data)
      at (0,\dataY) {Data Tier --- Firebase \\
        {\footnotesize\mdseries\color{ink} Cloud Firestore, Authentication, Cloud Functions}};

    % --- Element: Сум 1 (pres → app) -----------------------------------------
    \draw[hdarrow] (pres) -- node[right, font=\scriptsize] {REST API + Firebase SDK} (app);
    % --- Element: Сум 2 (app → data) -----------------------------------------
    \draw[hdarrow] (app)  -- node[right, font=\scriptsize] {Firestore Admin SDK} (data);

    % --- Element: Хажуугийн тэмдэглэгээ — клиент -----------------------------
    \node[font=\scriptsize\itshape\color{accent},  rotate=-8] (labelL) at (-6.5,\presY)  {клиент};
    % --- Element: Хажуугийн тэмдэглэгээ — логик -------------------------------
    \node[font=\scriptsize\itshape\color{accent2}, rotate=-8] (labelM) at (-6.5,\appY)   {логик};
    % --- Element: Хажуугийн тэмдэглэгээ — өгөгдөл -----------------------------
    \node[font=\scriptsize\itshape\color{accent3}, rotate=-8] (labelR) at (-6.5,\dataY)  {өгөгдөл};
  \end{tikzpicture}
\end{frame}

%###############################################################################
% SLIDE 6 — ТЕХНОЛОГИЙН СТЕК (хүснэгт)
%###############################################################################
\begin{frame}{Бүлэг 1: Технологийн стек}
  \vspace{-0.1cm}
  \centering
  \renewcommand{\arraystretch}{1.25}
  \begin{tikzpicture}
    \node[hdbox=paper, inner sep=8pt] {
      \begin{tabular}{|>{\columncolor{paper2}}p{2.4cm}|p{4.4cm}|p{6.5cm}|}
        \hline
        \rowcolor{paper2}\bfseries Давхарга & \bfseries Технологи & \bfseries Сонгосон шалтгаан\\
        \hline
        % --- Row: Flutter --------------------------------------------------
        Мобайл & Flutter 3.x + Dart & Нэг кодоос Android/iOS, Material Design 3\\
        \hline
        % --- Row: Firebase -------------------------------------------------
        Мобайл & Firebase Auth + Firestore & Real-time, аюулгүй нэвтрэлт\\
        \hline
        % --- Row: Cloud Functions ------------------------------------------
        Мобайл & Cloud Functions (Claude Haiku) & API key Flutter-т задрахгүй\\
        \hline
        % --- Row: Node.js --------------------------------------------------
        Бэкэнд & Node.js / Express & Event-driven, JSON API-д тохиромжтой\\
        \hline
        % --- Row: Groq -----------------------------------------------------
        Бэкэнд & Groq Llama 3.1-8b (нөөц AI) & LPU дэд бүтэц, маш хурдан\\
        \hline
        % --- Row: Docker ---------------------------------------------------
        Үүлэн & Docker & Орчны тогтвор, нээлттэй эх\\
        \hline
        % --- Row: AWS EC2 --------------------------------------------------
        Үүлэн & AWS EC2 (ap-southeast-1) & IaaS, бүрэн хяналт, Docker дэмжлэг\\
        \hline
        % --- Row: AWS ECR --------------------------------------------------
        Үүлэн & AWS ECR & Контейнер registry, IAM аюулгүй\\
        \hline
        % --- Row: GitHub Actions -------------------------------------------
        Үүлэн & GitHub Actions & CI/CD, git workflow-той нэгдмэл\\
        \hline
      \end{tabular}
    };
  \end{tikzpicture}
\end{frame}

%###############################################################################
% §3 — БҮЛЭГ 2 (хэсэг 3)
\section[Бүлэг 2]{Бүлэг 2: Програм хангамжийн хөгжүүлэлт}
%###############################################################################
% SLIDE 7 — БҮЛЭГ 2 НҮҮР
%###############################################################################
\begin{frame}[plain]
  \begin{tikzpicture}[remember picture, overlay]
    \fill[paper] (current page.north west) rectangle (current page.south east);
  \end{tikzpicture}
  \vspace{2cm}
  \centering
  \begin{tikzpicture}
    % --- Element: Гол гарчиг -----------------------------------------------
    \node[font=\Huge\bfseries\color{titlebg}] (t) {Бүлэг 2};
    % --- Element: Дэд гарчиг -----------------------------------------------
    \node[font=\LARGE\color{accent}, below=0.3cm of t] (s) {Програм хангамжийн хөгжүүлэлт};
    % --- Element: Доорх шугам ----------------------------------------------
    \draw[hd, draw=accent, line width=1.5pt]
      ($(s.south west)+(0.5,-0.2)$) -- ($(s.south east)+(-0.5,-0.2)$);
    % --- Element: Жижиг тайлбар --------------------------------------------
    \node[font=\small\itshape\color{ink!70}, below=1cm of s] (sub)
      {шаардлага · загварчлал · UI/UX · өгөгдлийн сан};
  \end{tikzpicture}
\end{frame}

%###############################################################################
% SLIDE 8 — USE CASE DIAGRAM
% Element-ууд:
%   sysBox  — системийн хүрээ (rectangle)
%   sysLbl  — системийн нэр
%   userL/userR — хоёр актор (Хэрэглэгч, Admin)
%   uc1..uc6 — use case-ууд
%   actorLine-* — актороос UC рүү шугам
%   includeRel — <<include>> dashed
%###############################################################################
\begin{frame}{Бүлэг 2.1: Use Case Diagram}
  \centering
  \begin{tikzpicture}[font=\scriptsize]
    % --- Element: Системийн хүрээ -------------------------------------------
    \draw[hd, draw=ink, line width=0.8pt, rounded corners=8pt]
      (-2.5,-3.2) rectangle (5.5,3.2);
    % --- Element: Системийн нэр ---------------------------------------------
    \node[font=\scriptsize\bfseries] (sysLbl) at (1.5,3.0) {Цаг захиалгын систем};

    % --- Element: Хэрэглэгч актор (зүүн) ------------------------------------
    \draw[hd, draw=ink] (-4.5,1.5) circle (0.25cm);   % толгой
    \draw[hd, draw=ink] (-4.5,1.25) -- (-4.5,0.5);    % бие
    \draw[hd, draw=ink] (-4.9,0.95) -- (-4.5,1.1) -- (-4.1,0.95);  % гар
    \draw[hd, draw=ink] (-4.5,0.5) -- (-4.8,-0.1);    % зүүн хөл
    \draw[hd, draw=ink] (-4.5,0.5) -- (-4.2,-0.1);    % баруун хөл
    \node[font=\scriptsize\bfseries, below] (userL) at (-4.5,-0.3) {Хэрэглэгч};

    % --- Element: Admin актор (баруун) -------------------------------------
    \draw[hd, draw=ink] (7,1.5) circle (0.25cm);
    \draw[hd, draw=ink] (7,1.25) -- (7,0.5);
    \draw[hd, draw=ink] (6.6,0.95) -- (7,1.1) -- (7.4,0.95);
    \draw[hd, draw=ink] (7,0.5) -- (6.7,-0.1);
    \draw[hd, draw=ink] (7,0.5) -- (7.3,-0.1);
    \node[font=\scriptsize\bfseries, below] (userR) at (7,-0.3) {Admin};

    % --- Element: Use Case 1 — Email бүртгэл --------------------------------
    \node[hdellipse=paper, minimum width=3.5cm] (uc1) at (1.5, 2.2)  {Email бүртгэл / нэвтрэх};
    % --- Element: Use Case 2 — Цагийн хуваарь ------------------------------
    \node[hdellipse=paper, minimum width=3.5cm] (uc2) at (1.5, 1.2)  {Цагийн хуваарь харах};
    % --- Element: Use Case 3 — Цаг захиалах --------------------------------
    \node[hdellipse=paper, minimum width=3.5cm] (uc3) at (1.5, 0.2)  {Цаг захиалах};
    % --- Element: Use Case 4 — Захиалга цуцлах ------------------------------
    \node[hdellipse=paper, minimum width=3.5cm] (uc4) at (1.5,-0.8)  {Захиалга цуцлах};
    % --- Element: Use Case 5 — AI чатбот -----------------------------------
    \node[hdellipse=paper, minimum width=3.5cm] (uc5) at (1.5,-1.8)  {AI чатбот харилцах};
    % --- Element: Use Case 6 — Тогтмол захиалга (Admin) ---------------------
    \node[hdellipse=paper2, minimum width=3.5cm] (uc6) at (1.5,-2.6) {Тогтмол захиалга удирдах};

    % --- Element: Хэрэглэгч → бүх UC --------------------------------------
    \foreach \u in {uc1,uc2,uc3,uc4,uc5}{
      \draw[hd, draw=ink, line width=0.5pt] (-4.25,0.95) -- (\u.west);
    }
    % --- Element: Admin → UC6 ------------------------------------------------
    \draw[hd, draw=ink, line width=0.5pt] (6.75,0.95) -- (uc6.east);
    % --- Element: <<include>> dashed харилцаа ------------------------------
    \draw[hddash, -{Latex[length=2mm]}] (uc3.south west) to[bend right=25]
      node[midway, sloped, above, font=\tiny\itshape] {<<include>>} (uc2.south west);
  \end{tikzpicture}
\end{frame}

%###############################################################################
% SLIDE 9 — ACTIVITY DIAGRAM
% Element-ууд:
%   s, e         — эхлэл/төгсгөл цэг
%   login..done  — алхмууд
%   avail        — diamond шийдвэр
%   err          — алдаа алхам (улаан)
%   arr-*        — алхам бүрийн сум
%###############################################################################
\begin{frame}{Бүлэг 2.1: Цаг захиалах Activity Diagram}
  \centering
  \begin{tikzpicture}[font=\scriptsize, node distance=0.55cm]
    % --- Element: Эхлэл цэг -------------------------------------------------
    \node[hd, draw=ink, fill=ink, circle, inner sep=2.5pt] (s) {};
    % --- Element: Алхам 1 — Нэвтрэлт ----------------------------------------
    \node[hdbox=paper, below=of s, minimum width=4cm]    (login)  {Email/нууц үгээр нэвтрэх};
    % --- Element: Алхам 2 — Огноо сонгох ------------------------------------
    \node[hdbox=paper, below=of login, minimum width=4cm](date)   {Огноо сонгох};
    % --- Element: Алхам 3 — Цаг сонгох --------------------------------------
    \node[hdbox=paper, below=of date, minimum width=4cm] (slot)   {Цагийн үе сонгох};
    % --- Element: Шийдвэр — боломжтой эсэх ----------------------------------
    \node[hddiamond=paper2, below=of slot, minimum width=2.5cm](avail) {Цаг боломжтой?};
    % --- Element: Алдааны алхам ---------------------------------------------
    \node[hdbox=paper, right=2cm of avail, minimum width=3cm, fill=accent!15] (err)
      {Алдаа --- дахин сонгох};
    % --- Element: Алхам 4 — Захиалга үүсгэх ---------------------------------
    \node[hdbox=paper, below=of avail, minimum width=4cm] (create) {Захиалга үүсгэх (Firestore)};
    % --- Element: Алхам 5 — DBK код ------------------------------------------
    \node[hdbox=paper, below=of create, minimum width=4cm](code)   {DBK-xxxxx код};
    % --- Element: Алхам 6 — Мэдэгдэл -----------------------------------------
    \node[hdbox=paper, below=of code, minimum width=4cm] (notif)  {1 цагийн өмнөх мэдэгдэл};
    % --- Element: Алхам 7 — Баталгаажуулалт ----------------------------------
    \node[hdbox=paper2, below=of notif, minimum width=4cm](done)  {Баталгаажуулалт};
    % --- Element: Төгсгөлийн цэг --------------------------------------------
    \node[hd, draw=ink, circle, double, double distance=1.5pt, inner sep=2pt, below=of done] (e) {};

    % --- Element: Сумнууд — гол урсгал --------------------------------------
    \draw[hdarrow] (s)     -- (login);
    \draw[hdarrow] (login) -- (date);
    \draw[hdarrow] (date)  -- (slot);
    \draw[hdarrow] (slot)  -- (avail);
    \draw[hdarrow] (avail) -- node[right, font=\tiny] {Тийм} (create);
    \draw[hdarrow] (create)-- (code);
    \draw[hdarrow] (code)  -- (notif);
    \draw[hdarrow] (notif) -- (done);
    \draw[hdarrow] (done)  -- (e);
    % --- Element: Шийдвэрээс алдаа руу --------------------------------------
    \draw[hdarrow] (avail) -- node[above, font=\tiny] {Үгүй} (err);
    % --- Element: Алдааны буцах урсгал --------------------------------------
    \draw[hdarrow] (err.north)|- ([xshift=2cm]slot.east) -- (slot);
  \end{tikzpicture}
\end{frame}

%###############################################################################
% SLIDE 10 — EMAIL БАТАЛГААЖУУЛАЛТЫН SEQUENCE DIAGRAM
% Element-ууд:
%   u, a, fb, b, e   — 5 lifeline толгой
%   life-*           — lifeline-ийн dashed шугам
%   m1..m8           — message-ууд (зүүнээс баруунаас солбицсон)
%###############################################################################
\begin{frame}{Бүлэг 2.2: Email Баталгаажуулалтын Систем}
  \centering
  \begin{tikzpicture}[font=\scriptsize, node distance=1.5cm,
                      every node/.append style={align=center}]
    % --- Element: Lifeline толгойнууд ---------------------------------------
    \node[hdbox=paper, minimum width=2.5cm, minimum height=1.2cm]              (u)  {\bfseries Хэрэглэгч};
    \node[hdbox=paper, minimum width=2.5cm, minimum height=1.2cm, right=of u]  (a)  {\bfseries Flutter App};
    \node[hdbox=paper, minimum width=2.5cm, minimum height=1.2cm, right=of a]  (fb) {\bfseries Firebase\\Auth};
    \node[hdbox=paper, minimum width=2.5cm, minimum height=1.2cm, right=of fb] (b)  {\bfseries Node.js\\EC2};
    \node[hdbox=paper, minimum width=2.5cm, minimum height=1.2cm, right=of b]  (e)  {\bfseries Gmail\\SMTP};

    % --- Element: Lifeline-ийн dashed шугам ---------------------------------
    \foreach \n in {u,a,fb,b,e}{
      \draw[hddash] (\n.south) -- ++(0,-5);
    }

    % --- Element: Message 1 — email/нууц үг (хэрэглэгч → апп) ----------------
    \draw[hdarrow] ($(u.south)+(0,-0.7)$)  -- node[above, font=\tiny] {1. email, нууц үг} ($(a.south)+(0,-0.7)$);
    % --- Element: Message 2 — createUser ------------------------------------
    \draw[hdarrow] ($(a.south)+(0,-1.3)$)  -- node[above, font=\tiny] {2. createUser} ($(fb.south)+(0,-1.3)$);
    % --- Element: Message 3 — UID буцах -------------------------------------
    \draw[hdarrow, draw=accent2] ($(fb.south)+(0,-1.9)$) -- node[above, font=\tiny] {3. UID} ($(a.south)+(0,-1.9)$);
    % --- Element: Message 4 — POST /send-code -------------------------------
    \draw[hdarrow] ($(a.south)+(0,-2.5)$)  -- node[above, font=\tiny] {4. POST /send-code} ($(b.south)+(0,-2.5)$);
    % --- Element: Message 5 — 6-digit код илгээх ----------------------------
    \draw[hdarrow] ($(b.south)+(0,-3.1)$)  -- node[above, font=\tiny] {5. 6-digit code} ($(e.south)+(0,-3.1)$);
    % --- Element: Message 6 — OK хариу --------------------------------------
    \draw[hdarrow, draw=accent2] ($(e.south)+(0,-3.7)$) -- node[above, font=\tiny] {6. OK} ($(b.south)+(0,-3.7)$);
    % --- Element: Message 7 — 200 OK ----------------------------------------
    \draw[hdarrow, draw=accent2] ($(b.south)+(0,-4.3)$) -- node[above, font=\tiny] {7. 200 OK} ($(a.south)+(0,-4.3)$);
    % --- Element: Message 8 — Код оруулах -----------------------------------
    \draw[hdarrow, draw=accent2] ($(a.south)+(0,-4.9)$) -- node[above, font=\tiny] {8. Код оруул} ($(u.south)+(0,-4.9)$);
  \end{tikzpicture}
\end{frame}

%###############################################################################
% SLIDE 11 — ДАВХЦАЛ ШАЛГАХ ЛОГИК
% Element-ууд: codeBox (Firestore Query), notesBox (тайлбар)
%###############################################################################
\begin{frame}{Бүлэг 2.2: Давхцал шалгах логик --- Firestore Query}
  \centering
  \begin{tikzpicture}[font=\footnotesize]
    % --- Element: Код блок ---------------------------------------------------
    \node[hdbox=paper, text width=12cm, align=left, font=\ttfamily\footnotesize, inner sep=10pt] (codeBox) {%
      \textbf{// Захиалга үүсгэхээс өмнө давхардал шалгах}\\[2pt]
      final snap = await db.collection('bookings')\\
      \hspace*{1em}.where('venueId', isEqualTo: venueId)\\
      \hspace*{1em}.where('date',    isEqualTo: date)\\
      \hspace*{1em}.where('timeSlot',isEqualTo: slot)\\
      \hspace*{1em}.where('status',  whereIn: ['upcoming','active'])\\
      \hspace*{1em}.get();\\[4pt]
      \textbf{if} (snap.docs.isNotEmpty) \{\\
      \hspace*{1em}\textbf{throw} 'Уг цаг авагдсан байна';\\
      \}};
    % --- Element: Тайлбар блок ----------------------------------------------
    \node[hdbox=paper2, text width=10cm, font=\scriptsize, anchor=north] (notesBox)
      at (0,-3.2) {\textbf{Compound index} ашигласан тул query 50ms-д дуусдаг.\\
        Хоёр хэрэглэгч нэгэн зэрэг захиалбал \textbf{atomic transaction}-аар нэг нь л баталгаажна.};
  \end{tikzpicture}
\end{frame}

%###############################################################################
% SLIDE 12 — UI/UX ДЭЛГЭЦҮҮД
% Element-ууд: phone-1..4 (4 утасны хүрээ), descBox
%###############################################################################
\begin{frame}{Бүлэг 2.2: Мобайл аппликейшн --- UI/UX}
  \centering
  \begin{tikzpicture}[font=\scriptsize]
    % --- 4 утасны mock-up -- одоохондоо хүрээ зурдаг макрог давтсан ---------
    \foreach \i/\title in {%
      0/{Нүүр},%
      1/{Цаг захиалах},%
      2/{Захиалга баталгаажуулах},%
      3/{Миний захиалгууд}}{

        \node[hdbox=paper, minimum width=3cm, minimum height=4.5cm,
              text width=2.6cm, font=\scriptsize, inner sep=4pt]
          (s\i) at ({-6+\i*4}, 0.5) {%
            \vspace{0.2cm}
            \begin{tikzpicture}
              % Element: Утасны бие (дэвсгэр)
              \draw[hd, draw=ink!50, fill=paper2] (0,0) rectangle (2,3.5);
              % Element: Толгой бар
              \draw[hd, draw=accent3, fill=accent3!20] (0.1,3.0) rectangle (1.9,3.4);
              % Element: 1-р мөр
              \draw[hd, draw=ink!60] (0.2,2.4) rectangle (1.8,2.8);
              % Element: 2-р мөр
              \draw[hd, draw=ink!60] (0.2,1.7) rectangle (1.8,2.1);
              % Element: 3-р мөр
              \draw[hd, draw=ink!60] (0.2,1.0) rectangle (1.8,1.4);
              % Element: Үндсэн товч
              \draw[hd, draw=accent, fill=accent!30] (0.2,0.2) rectangle (1.8,0.7);
            \end{tikzpicture}};
        % --- Element: дэлгэцийн нэр --------------------------------------
        \node[font=\scriptsize\bfseries, below=0.05cm of s\i] {\title};
      }

    % --- Element: Доод тайлбар ----------------------------------------------
    \node[hdbox=paper2, text width=14cm, font=\scriptsize\itshape, anchor=north] (descBox)
      at (0,-3.5)
      {Material Design 3 + дулаан говийн өнгөн схем (бежий \#EDE7DA, наранжин \#D4700F)};
  \end{tikzpicture}
\end{frame}

%###############################################################################
% SLIDE 13 — AI ТУСЛАХ (Хоёр горим)
% Element-ууд: app, cf, ec2, cl, gr — 5 box; arr-* сумнууд
%###############################################################################
\begin{frame}{Бүлэг 2.2: AI Туслах --- Хоёр горим}
  \centering
  \begin{tikzpicture}[font=\footnotesize, node distance=1cm]
    % --- Element: Flutter App box -------------------------------------------
    \node[hdbox=paper, minimum width=3.5cm, minimum height=1.2cm] (app) at (-5,1.5)
      {\bfseries Flutter App};
    % --- Element: Cloud Functions box (үндсэн горим) -----------------------
    \node[hdbox=paper, minimum width=4cm, minimum height=1.2cm, fill=accent!15] (cf)
      at (0,2.5) {\bfseries Firebase\\Cloud Functions};
    % --- Element: Node.js / EC2 box (нөөц горим) ----------------------------
    \node[hdbox=paper, minimum width=4cm, minimum height=1.2cm, fill=accent2!15] (ec2)
      at (0,0.5) {\bfseries Node.js / EC2};
    % --- Element: Claude Haiku box ------------------------------------------
    \node[hdbox=paper, minimum width=3.2cm, minimum height=1.2cm] (cl) at (5,2.5)
      {\bfseries Claude Haiku 4.5\\(үндсэн)};
    % --- Element: Groq Llama box --------------------------------------------
    \node[hdbox=paper, minimum width=3.2cm, minimum height=1.2cm] (gr) at (5,0.5)
      {\bfseries Groq Llama 3.1\\(нөөц)};

    % --- Element: Сум 1 (App → CF) ------------------------------------------
    \draw[hdarrow] (app) -- node[above, font=\tiny] {Cloud Function} (cf);
    % --- Element: Сум 2 (App → EC2 fallback) --------------------------------
    \draw[hdarrow, draw=accent] (app) to[bend right=20]
      node[below, font=\tiny\itshape] {fallback if timeout} (ec2);
    % --- Element: Сум 3 (CF → Claude) ---------------------------------------
    \draw[hdarrow] (cf)  -- (cl);
    % --- Element: Сум 4 (EC2 → Groq) -----------------------------------------
    \draw[hdarrow] (ec2) -- (gr);

    % --- Element: Доод тайлбар ----------------------------------------------
    \node[hdbox=paper2, text width=14cm, font=\scriptsize, anchor=north] (descBox)
      at (0,-1.5){
        \textbf{Үндсэн горим} нь Claude Haiku-аар дуудагдаж $<$ 1 секундэд хариулдаг.\\
        \textbf{Нөөц горим} (Groq Llama) нь Claude алдаа гарсан тохиолдолд автоматаар идэвхждэг.};
  \end{tikzpicture}
\end{frame}

%###############################################################################
% SLIDE 14 — FIRESTORE ERD
% Element-ууд: 4 collection (u, b, f, e); 2 relation (rel-userId, rel-uid)
%###############################################################################
\begin{frame}{Бүлэг 2.4: Firebase Firestore --- Өгөгдлийн схем}
  \centering
  \begin{tikzpicture}[font=\scriptsize, every node/.style={align=left}]
    % --- Element: users collection ------------------------------------------
    \node[hdbox=paper, text width=4cm, font=\scriptsize, inner sep=6pt] (u) at (-4.5, 2) {%
      {\bfseries\color{accent3} users}\\[2pt]
      uid (doc ID)\\
      name : string\\
      email : string\\
      emailVerified : bool\\
      createdAt : timestamp};
    % --- Element: bookings collection ---------------------------------------
    \node[hdbox=paper, text width=4cm, font=\scriptsize, inner sep=6pt] (b) at (4.5, 2) {%
      {\bfseries\color{accent3} bookings}\\[2pt]
      code : DBK-xxxxx\\
      userId : string\\
      venueId : string\\
      date / dateKey\\
      timeSlot / timeSlotEnd\\
      status : upcoming\textbar cancelled};
    % --- Element: fixed_bookings collection ----------------------------------
    \node[hdbox=paper2, text width=4cm, font=\scriptsize, inner sep=6pt] (f) at (-4.5,-1.8) {%
      {\bfseries\color{accent2} fixed\_bookings}\\[2pt]
      venueId : string\\
      organizationName\\
      dayOfWeek : 0--6\\
      timeSlot : string\\
      isActive : bool};
    % --- Element: email_verifications collection ----------------------------
    \node[hdbox=paper2, text width=4cm, font=\scriptsize, inner sep=6pt] (e) at (4.5,-1.8) {%
      {\bfseries\color{accent} email\_verifications}\\[2pt]
      uid (doc ID)\\
      code : 6 оронтой\\
      email : string\\
      expiresAt : timestamp};

    % --- Element: Relation users.uid → bookings.userId ---------------------
    \draw[hddash, draw=ink] (u.east)  -- node[above, font=\tiny] {userId} (b.west);
    % --- Element: Relation users.uid → email_verifications.uid -------------
    \draw[hddash, draw=ink] (u.south) -- node[right, font=\tiny] {uid} (e.north);
  \end{tikzpicture}
\end{frame}

%###############################################################################
% §4 — БҮЛЭГ 3 (хэсэг 4)
\section[Бүлэг 3]{Бүлэг 3: Үүлэн орчны хэрэгжүүлэлт}
%###############################################################################
% SLIDE 15 — БҮЛЭГ 3 НҮҮР
%###############################################################################
\begin{frame}[plain]
  \begin{tikzpicture}[remember picture, overlay]
    \fill[paper] (current page.north west) rectangle (current page.south east);
  \end{tikzpicture}
  \vspace{2cm}
  \centering
  \begin{tikzpicture}
    % --- Element: Гол гарчиг -----------------------------------------------
    \node[font=\Huge\bfseries\color{titlebg}] (t) {Бүлэг 3};
    % --- Element: Дэд гарчиг ------------------------------------------------
    \node[font=\LARGE\color{accent}, below=0.3cm of t] (s) {Үүлэн орчны хэрэгжүүлэлт};
    % --- Element: Шугам -----------------------------------------------------
    \draw[hd, draw=accent, line width=1.5pt]
      ($(s.south west)+(0.5,-0.2)$) -- ($(s.south east)+(-0.5,-0.2)$);
    % --- Element: Жижиг тайлбар --------------------------------------------
    \node[font=\small\itshape\color{ink!70}, below=1cm of s]
      {Docker · AWS EC2 · GitHub Actions CI/CD · туршилт};
  \end{tikzpicture}
\end{frame}

%###############################################################################
% SLIDE 16 — DEPLOYMENT OVERVIEW
% Element-ууд: gh, act, ecr, ec2, ph, ps — 6 box; сум 5 ш
%###############################################################################
\begin{frame}{Бүлэг 3.1: Deployment overview}
  \centering
  \begin{tikzpicture}[font=\footnotesize]
    % --- Element: GitHub Repo box -------------------------------------------
    \node[hdbox=paper, minimum width=3cm, minimum height=1.4cm] (gh)
      at (-5.5, 2)  {\bfseries GitHub\\Repository};
    % --- Element: GitHub Actions box ----------------------------------------
    \node[hdbox=paper, minimum width=3cm, minimum height=1.4cm] (act)
      at (-5.5,-1) {\bfseries GitHub\\Actions CI/CD};
    % --- Element: AWS ECR box ------------------------------------------------
    \node[hdbox=paper, minimum width=3cm, minimum height=1.4cm, fill=accent!15] (ecr)
      at ( 0,-1) {\bfseries AWS ECR\\Docker Registry};
    % --- Element: AWS EC2 box ------------------------------------------------
    \node[hdbox=paper, minimum width=3cm, minimum height=1.4cm, fill=accent2!15] (ec2)
      at ( 0, 2) {\bfseries AWS EC2\\Docker контейнер};
    % --- Element: Flutter App box -------------------------------------------
    \node[hdbox=paper, minimum width=3cm, minimum height=1.4cm] (ph)
      at ( 5.5, 2) {\bfseries Flutter App\\Android / iOS};
    % --- Element: Play Store box --------------------------------------------
    \node[hdbox=paper, minimum width=3cm, minimum height=1.4cm, fill=accent3!15] (ps)
      at ( 5.5,-1) {\bfseries Google Play Store};

    % --- Element: Сум 1 (GitHub → Actions: push) ----------------------------
    \draw[hdarrow] (gh)  -- node[left, font=\tiny] {push}         (act);
    % --- Element: Сум 2 (Actions → ECR: docker push) ------------------------
    \draw[hdarrow] (act) -- node[above, font=\tiny] {docker push} (ecr);
    % --- Element: Сум 3 (ECR → EC2: pull) -----------------------------------
    \draw[hdarrow] (ecr) -- node[left, font=\tiny] {pull}         (ec2);
    % --- Element: Сум 4 (EC2 → Flutter: REST API) ---------------------------
    \draw[hdarrow] (ec2) -- node[above, font=\tiny] {REST API}    (ph);
    % --- Element: Сум 5 (Play Store → Flutter: download) --------------------
    \draw[hdarrow] (ps)  -- node[right, font=\tiny] {download}    (ph);

    % --- Element: Доод тайлбар ----------------------------------------------
    \node[font=\scriptsize\itshape\color{accent}, anchor=south]
      at (0,-2.7) {3--4 минутын автомат байршуулалт};
  \end{tikzpicture}
\end{frame}

%###############################################################################
% SLIDE 17 — CI/CD PIPELINE
% Element-ууд: p1..p7 алхмууд, fail (улаан box), сум-ууд
%###############################################################################
\begin{frame}{Бүлэг 3.1: CI/CD Pipeline --- Docker + AWS EC2}
  \centering
  \begin{tikzpicture}[font=\scriptsize, node distance=0.55cm]
    % --- Element: Алхам 1 — git push ----------------------------------------
    \node[hdbox=paper, minimum width=3.5cm]              (p1) {git push (\textit{backend/})};
    % --- Element: Алхам 2 — GH Actions trigger -------------------------------
    \node[hdbox=paper, minimum width=3.5cm, below=of p1] (p2) {GH Actions trigger};
    % --- Element: Алхам 3 — Docker build -------------------------------------
    \node[hdbox=paper, minimum width=3.5cm, below=of p2] (p3) {Docker build};
    % --- Element: Шийдвэр — Build OK эсэх -----------------------------------
    \node[hddiamond=paper2, below=of p3, minimum width=2.5cm] (p4) {Build OK?};
    % --- Element: Алхам 5 — ECR push ----------------------------------------
    \node[hdbox=paper, minimum width=3.5cm, below=of p4] (p5) {ECR push};
    % --- Element: Алхам 6 — SSH to EC2 ---------------------------------------
    \node[hdbox=paper, minimum width=3.5cm, below=of p5] (p6) {SSH to EC2};
    % --- Element: Алхам 7 — docker pull / run --------------------------------
    \node[hdbox=paper, minimum width=3.5cm, below=of p6] (p7) {docker pull / run};
    % --- Element: Алдааны замбраа --------------------------------------------
    \node[hdbox=paper, fill=accent!15, right=2.8cm of p4, minimum width=3cm] (fail)
      {Workflow fail};

    % --- Element: Сумнууд (гол урсгал) --------------------------------------
    \foreach \a/\b in {p1/p2,p2/p3,p3/p4,p4/p5,p5/p6,p6/p7}{
      \draw[hdarrow] (\a) -- (\b);
    }
    % --- Element: Үгүй сум — fail руу ---------------------------------------
    \draw[hdarrow] (p4) -- node[above, font=\tiny] {Үгүй} (fail);
    % --- Element: "Тийм" текст ----------------------------------------------
    \node[font=\tiny] at ($(p4.south)+(0.6,-0.1)$) {Тийм};

    % --- Element: Тайлбарын box ----------------------------------------------
    \node[hdbox=paper2, font=\scriptsize, text width=4cm, align=left, right=1.2cm of p7] (notes)
      {\textbf{Үе шат}: 6 алхамтай\\
       \textbf{Хугацаа}: 3--4 минут\\
       \textbf{Зэрэгцэх}: бүтэн автомат\\
       \textbf{Хяналт}: GitHub UI};
  \end{tikzpicture}
\end{frame}

%###############################################################################
% SLIDE 18 — ХАРИУ ЦАГИЙН ХЭМЖИЛТ (xbar диаграмм)
%###############################################################################
\begin{frame}{Бүлэг 3.2: Хариу цагийн хэмжилт ($n=10$, 4G LTE)}
  \centering
  \begin{tikzpicture}[font=\small]
    % --- Element: График ----------------------------------------------------
    \node[hdbox=paper, inner sep=8pt] {
      \begin{tikzpicture}[font=\footnotesize]
        \begin{axis}[
          xbar, width=11cm, height=5.8cm, bar width=8pt,
          xmin=0, xmax=5.6, xlabel={Хариу цаг (секунд)},
          ytick={1,2,3,4,5},
          yticklabels={{AI чатбот},{Захиалга үүсгэх},{EC2 API},{Firestore},{Апп ачаалах}},
          y tick label style={font=\footnotesize},
          nodes near coords, nodes near coords align=horizontal,
          every node near coord/.append style={font=\tiny},
          enlarge y limits=0.18,
          tick style={draw=ink}, axis line style={draw=ink},
          legend style={at={(0.98,0.05)}, anchor=south east, font=\scriptsize, draw=ink},
          legend cell align=left]
          % Element: Хэмжсэн утгууд
          \addplot[fill=accent3!40, draw=accent3] coordinates
            {(0.9,1)(1.1,2)(0.8,3)(0.6,4)(2.3,5)};
          % Element: Шаардлагын хязгаар
          \addplot[fill=accent!20, draw=accent, dashed] coordinates
            {(5.0,1)(2.0,2)(2.0,3)(2.0,4)(3.0,5)};
          \legend{Хэмжсэн (с), Шаардлага (с)}
        \end{axis}
      \end{tikzpicture}};
    % --- Element: Доод тайлбар ----------------------------------------------
    \node[hdbox=paper2, text width=14cm, font=\scriptsize\itshape, anchor=north]
      at (0,-3.5){Бүх үзүүлэлт шаардлагын хязгаарт нийцсэн.};
  \end{tikzpicture}
\end{frame}

%###############################################################################
% SLIDE 19 — ХЭРЭГЖИЛТИЙН ҮР ДҮН (хүснэгт)
%###############################################################################
\begin{frame}{Бүлэг 3.2: Хэрэгжилтийн үр дүн}
  \centering
  \renewcommand{\arraystretch}{1.2}
  \begin{tikzpicture}
    \node[hdbox=paper, inner sep=8pt] {
      \begin{tabular}{|p{4.5cm}|p{6cm}|p{2.6cm}|}
        \hline
        \rowcolor{paper2}\bfseries Шалгах зүйл & \bfseries Үзүүлэлт & \bfseries Үр дүн\\
        \hline
        % --- Row: Flutter мобайл апп -----------------------------------------
        Flutter мобайл апп & 8 дэлгэц, iOS / Android & {\color{accent2}\textbf{DONE}}\\
        \hline
        % --- Row: Node.js REST API -------------------------------------------
        Node.js REST API & 8 endpoint, production EC2 & {\color{accent2}\textbf{DONE}}\\
        \hline
        % --- Row: Firebase нэгтгэл --------------------------------------------
        Firebase нэгтгэл & Firestore + Auth OTP + локал мэдэгдэл & {\color{accent2}\textbf{DONE}}\\
        \hline
        % --- Row: AI Chat ----------------------------------------------------
        AI Chat туслах & Claude Haiku + Groq Llama, $<$1 сек & {\color{accent2}\textbf{DONE}}\\
        \hline
        % --- Row: Давхцал шалгалт ---------------------------------------------
        Давхцал шалгалт & Compound index, atomic transaction & {\color{accent2}\textbf{DONE}}\\
        \hline
        % --- Row: CI/CD ------------------------------------------------------
        CI/CD Pipeline & GitHub Actions автомат деплой & {\color{accent2}\textbf{DONE}}\\
        \hline
        % --- Row: Туршилт ----------------------------------------------------
        Туршилт (нэгтгэсэн) & 9 функциональ + 5 гүйцэтгэлийн алхам & {\color{accent2}\textbf{DONE}}\\
        \hline
        % --- Row: Хэрэглэхэд хялбар -------------------------------------------
        Хэрэглэхэд хялбар & 4.3 / 5.0 оноо (Usability scale) & {\color{accent2}\textbf{DONE}}\\
        \hline
      \end{tabular}};
  \end{tikzpicture}
\end{frame}

%###############################################################################
% §5 — ДҮГНЭЛТ (хэсэг 5)
\section[Дүгнэлт]{Дүгнэлт}
%###############################################################################
% SLIDE 20 — ДҮГНЭЛТ (3 баган)
% Element-ууд: topBox, col-1 шинэлэг, col-2 ач холбогдол, col-3 цаашид
%###############################################################################
\begin{frame}{Дүгнэлт}
  \begin{tikzpicture}[font=\footnotesize, node distance=0.45cm]
    % --- Element: Дүгнэлтийн дүн ---------------------------------------------
    \node[hdbox=paper2, text width=14cm, align=left, font=\small, inner sep=10pt] (sumBox) {%
      Flutter (Presentation), Node.js (Application), Firebase (Data) давхаргуудыг \textbf{AWS EC2} дэд бүтэцэд суурилуулан нэгтгэсэн \textbf{3 давхаргат} ухаалаг захиалгын систем амжилттай загварчлагдаж, production орчинд хэрэгжлээ.};

    % --- Element: 1-р баган — Шинэлэг тал ----------------------------------
    \node[hdbox=paper, font=\scriptsize, text width=4.4cm, align=left, anchor=north west,
          inner sep=6pt] (col1)
      at (-7.2,-0.5) {\textbf{\color{accent3} Шинэлэг тал}\\
        \textbullet\ 3 давхаргат үүлэн архитектур\\
        \textbullet\ Давхар горимын AI чатбот\\
        \textbullet\ \texttt{fixed\_bookings} функц};

    % --- Element: 2-р баган — Ач холбогдол ---------------------------------
    \node[hdbox=paper, font=\scriptsize, text width=4.4cm, align=left, anchor=north,
          inner sep=6pt] (col2)
      at (0,-0.5) {\textbf{\color{accent2} Ач холбогдол}\\
        \textbullet\ Иргэдэд хялбар\\
        \textbullet\ Захиалгын ачаалал бууруулна\\
        \textbullet\ Орон нутгийн жишиг};

    % --- Element: 3-р баган — Цаашдын хөгжил -------------------------------
    \node[hdbox=paper, font=\scriptsize, text width=4.4cm, align=left, anchor=north east,
          inner sep=6pt] (col3)
      at (7.2,-0.5) {\textbf{\color{accent} Цаашдын хөгжил}\\
        \textbullet\ QPay/SocialPay төлбөр\\
        \textbullet\ HTTPS/SSL\\
        \textbullet\ Админ дэлгэц};
  \end{tikzpicture}
\end{frame}

%###############################################################################
% SLIDE 21 — АНХААРАЛ ТАВЬСАНД БАЯРЛАЛАА
%###############################################################################
\begin{frame}[plain]
  \begin{tikzpicture}[remember picture, overlay]
    \fill[paper] (current page.north west) rectangle (current page.south east);
  \end{tikzpicture}
  \vspace{1.5cm}
  \centering
  \begin{tikzpicture}
    % --- Element: Талархлын box ----------------------------------------------
    \node[hdbox=paper2, text width=12cm, font=\Huge\bfseries\color{titlebg}, inner sep=20pt] (t)
      {Анхаарал тавьсанд\\баярлалаа};
    % --- Element: Доорх шугам ------------------------------------------------
    \draw[hd, draw=accent, line width=1.5pt]
      ($(t.south west)+(1,-0.4)$) -- ($(t.south east)+(-1,-0.4)$);
  \end{tikzpicture}\\[0.6cm]
  % --- Element: Асуултын текст ----------------------------------------------
  {\Large\itshape\color{ink} Асуулт?}\\[1cm]
  % --- Element: Илтгэгчийн нэр ----------------------------------------------
  {\small\color{ink!70} М.Түвшинжаргал \quad B221940619 \quad 2026 он}
\end{frame}

\end{document}
